\documentclass[11pt]{article}
\usepackage{mathblatt}
\newcommand{\bereich}[1]{\par\addvspace{14pt}\noindent{\large\bfseries #1}\par\nobreak\vspace{4pt}}
\newcommand{\bl}[1]{\par\noindent #1\par\vspace{3pt}}
\newcommand{\blrand}[1]{\par\noindent{\color{mbgrau}#1}\par\vspace{3pt}}
\newcommand{\blhier}[1]{\par\noindent\llap{$\blacktriangleright$\hspace{0.5em}}\textbf{#1}\par\vspace{3pt}}
\newcommand{\stufe}[1]{\par\noindent\hspace*{1.5em}{\small #1}\par\vspace{2pt}}
\newcommand{\pkt}{\textperiodcentered{}}
\begin{document}

\pfheftstil
\pflaufkopf[Abitur GK \pkt{} Prüfung]{Kurvenuntersuchung}{AA60}\pfseitevon
\pfheftkopf{Kurvenuntersuchung}{}
\bl{Grenzverhalten und Symmetrie}
\bl{Nullstellen}
\bl{Extrempunkte}
\blhier{Monotonie}
\stufe{Monotonie begründen}
\stufe{Graph von f$'$ lesen}
\stufe{Aussagen über f und f$'$ beurteilen}
\bl{Wendepunkte}
\blrand{Gleichungen und Ungleichungen}
\blrand{Funktion aufstellen}
\blrand{Verschieben, spiegeln, strecken}
\blrand{Graph im Sachzusammenhang}

\par\addvspace{14pt}\noindent Weiter: Ableitung + Tangente\par
\end{document}